\section{洪熙至天顺的补充编年节点（一）}
这一组条目把此前尚未设导航入口的年度材料接回本章脉络。它们不是把同年材料机械等同为因果，而是在可核查的时间位置上，分别保留行动、行政记录与社会经验之间的距离。

\subsection{1434年前后的材料线索}

\eventanchor{ML-1434-001}{明·1434（宣德九年）}
\paragraph{1434（宣德九年）：宝藏之旅：龚震出版《西洋繁国志》} 这一节点应放回本章已经展开的权力、财政、地方社会与边地关系中理解。账本把它出现的条件概括为：继承、官僚协调或皇权运作的压力使问题进入朝廷程序。 随后可追踪的变化是：它影响后续人事与政策议程；动机和派系标签不能由单一叙事裁定。 可由《宣宗实录》宣德九年条；《明史·本纪》及相关列传；诏令、奏疏或会典版本 互相核对；不过，译写候选以编年材料定位；实录记载反映朝廷选择，崇祯期须另以奏疏、地方及敌对政权材料交叉。 因而这里标示的是材料能够支持的行动与制度位置，而不是对动机、地方执行或普通人经验的无保留复原。

\eventanchor{ML-1434-002}{明·1434（宣德九年）}
\paragraph{1434（宣德九年）：本年朝廷围绕继承、官僚协调和政令传达形成连续记录，可用于核对宣德至正统的中枢运作。} 这一节点应放回本章已经展开的权力、财政、地方社会与边地关系中理解。账本把它出现的条件概括为：继承、官僚协调或皇权运作的压力使问题进入朝廷程序。 随后可追踪的变化是：它影响后续人事与政策议程；动机和派系标签不能由单一叙事裁定。 可由《宣宗实录》宣德九年条；《明史·本纪》及相关列传；诏令、奏疏或会典版本 互相核对；不过，桥接条目只标示可回查的年度问题与材料链；实录有中央选择性，崇祯期尤其需要奏疏、地方、朝鲜及满文材料交叉。 因而这里标示的是材料能够支持的行动与制度位置，而不是对动机、地方执行或普通人经验的无保留复原。

\eventanchor{ML-1434-003}{明·1434（宣德九年）}
\paragraph{1434（宣德九年）：学校、科举和礼制的本年政策材料为社会流动和国家知识秩序提供可检索锚点。} 这一节点应放回本章已经展开的权力、财政、地方社会与边地关系中理解。账本把它出现的条件概括为：制度、教育或知识传播的需要推动了相关行动或文本形成。 随后可追踪的变化是：它提供制度和传播史的锚点，不能据此推断社会整体接受。 可由《宣宗实录》宣德九年条；《明史·选举志》或《艺文志》；文集、碑刻或书目材料 互相核对；不过，桥接条目只标示可回查的年度问题与材料链；实录有中央选择性，崇祯期尤其需要奏疏、地方、朝鲜及满文材料交叉。 因而这里标示的是材料能够支持的行动与制度位置，而不是对动机、地方执行或普通人经验的无保留复原。

\eventanchor{ML-1434-004}{明·1434（宣德九年）}
\paragraph{1434（宣德九年）：朝贡、册封和边地交涉的记录为本年多政权关系留档，礼仪语言与实际控制范围必须区分。} 这一节点应放回本章已经展开的权力、财政、地方社会与边地关系中理解。账本把它出现的条件概括为：朝贡名义、边境安全与港口贸易的利益使交涉持续发生。 随后可追踪的变化是：它为后续册封、互市或海防安排留下文书与跨政权比较线索。 可由《宣宗实录》宣德九年条；《明史·外国传》；《朝鲜王朝实录》、琉球《历代宝案》或海外档案 互相核对；不过，桥接条目只标示可回查的年度问题与材料链；实录有中央选择性，崇祯期尤其需要奏疏、地方、朝鲜及满文材料交叉。 因而这里标示的是材料能够支持的行动与制度位置，而不是对动机、地方执行或普通人经验的无保留复原。

\eventanchor{ML-1435-001}{明·1435（宣德十年）一月31日}
\paragraph{1435（宣德十年）一月31日：宣德皇帝驾崩，张皇后（洪熙）摄政为正统皇帝} 这一节点应放回本章已经展开的权力、财政、地方社会与边地关系中理解。账本把它出现的条件概括为：继承、官僚协调或皇权运作的压力使问题进入朝廷程序。 随后可追踪的变化是：它影响后续人事与政策议程；动机和派系标签不能由单一叙事裁定。 可由《宣宗实录》宣德十年条；《明史·本纪》及相关列传；诏令、奏疏或会典版本 互相核对；不过，译写候选以编年材料定位；实录记载反映朝廷选择，崇祯期须另以奏疏、地方及敌对政权材料交叉。 因而这里标示的是材料能够支持的行动与制度位置，而不是对动机、地方执行或普通人经验的无保留复原。

\eventanchor{ML-1435-002}{明·1435（宣德十年）}
\paragraph{1435（宣德十年）：华北平原及山东遭遇旱灾、瘟疫} 这一节点应放回本章已经展开的权力、财政、地方社会与边地关系中理解。账本把它出现的条件概括为：自然风险与地方报告促成朝廷或地方的应对记录。 随后可追踪的变化是：赈济、迁徙和损失的实际范围仍须以受灾地区材料分别核验。 可由《宣宗实录》宣德十年条；《明史·五行志》；受灾府县地方志与奏疏 互相核对；不过，译写候选以编年材料定位；实录记载反映朝廷选择，崇祯期须另以奏疏、地方及敌对政权材料交叉。 因而这里标示的是材料能够支持的行动与制度位置，而不是对动机、地方执行或普通人经验的无保留复原。

\eventanchor{ML-1435-003}{明·1435（宣德十年）}
\paragraph{1435（宣德十年）：学校、科举和礼制的本年政策材料为社会流动和国家知识秩序提供可检索锚点。} 这一节点应放回本章已经展开的权力、财政、地方社会与边地关系中理解。账本把它出现的条件概括为：制度、教育或知识传播的需要推动了相关行动或文本形成。 随后可追踪的变化是：它提供制度和传播史的锚点，不能据此推断社会整体接受。 可由《宣宗实录》宣德十年条；《明史·选举志》或《艺文志》；文集、碑刻或书目材料 互相核对；不过，桥接条目只标示可回查的年度问题与材料链；实录有中央选择性，崇祯期尤其需要奏疏、地方、朝鲜及满文材料交叉。 因而这里标示的是材料能够支持的行动与制度位置，而不是对动机、地方执行或普通人经验的无保留复原。

\eventanchor{ML-1436-001}{明·1436（正统元年）}
\paragraph{1436（正统元年）：宝藏航行：明朝禁止建造海船} 这一节点应放回本章已经展开的权力、财政、地方社会与边地关系中理解。账本把它出现的条件概括为：继承、官僚协调或皇权运作的压力使问题进入朝廷程序。 随后可追踪的变化是：它影响后续人事与政策议程；动机和派系标签不能由单一叙事裁定。 可由《英宗实录》正统元年条；《明史·本纪》及相关列传；诏令、奏疏或会典版本 互相核对；不过，译写候选以编年材料定位；实录记载反映朝廷选择，崇祯期须另以奏疏、地方及敌对政权材料交叉。 因而这里标示的是材料能够支持的行动与制度位置，而不是对动机、地方执行或普通人经验的无保留复原。

\eventanchor{ML-1436-002}{明·1436（正统元年）}
\paragraph{1436（正统元年）：宝藏之旅：费辛出版《星茶圣澜》} 这一节点应放回本章已经展开的权力、财政、地方社会与边地关系中理解。账本把它出现的条件概括为：继承、官僚协调或皇权运作的压力使问题进入朝廷程序。 随后可追踪的变化是：它影响后续人事与政策议程；动机和派系标签不能由单一叙事裁定。 可由《英宗实录》正统元年条；《明史·本纪》及相关列传；诏令、奏疏或会典版本 互相核对；不过，译写候选以编年材料定位；实录记载反映朝廷选择，崇祯期须另以奏疏、地方及敌对政权材料交叉。 因而这里标示的是材料能够支持的行动与制度位置，而不是对动机、地方执行或普通人经验的无保留复原。

\eventanchor{ML-1436-003}{明·1436（正统元年）}
\paragraph{1436（正统元年）：苏北、华北平原、山东等地遭受洪涝灾害} 这一节点应放回本章已经展开的权力、财政、地方社会与边地关系中理解。账本把它出现的条件概括为：自然风险与地方报告促成朝廷或地方的应对记录。 随后可追踪的变化是：赈济、迁徙和损失的实际范围仍须以受灾地区材料分别核验。 可由《英宗实录》正统元年条；《明史·五行志》；受灾府县地方志与奏疏 互相核对；不过，译写候选以编年材料定位；实录记载反映朝廷选择，崇祯期须另以奏疏、地方及敌对政权材料交叉。 因而这里标示的是材料能够支持的行动与制度位置，而不是对动机、地方执行或普通人经验的无保留复原。

\eventanchor{ML-1437-001}{明·1437（正统二年）}
\paragraph{1437（正统二年）：山西、陕西遭遇干旱} 这一节点应放回本章已经展开的权力、财政、地方社会与边地关系中理解。账本把它出现的条件概括为：自然风险与地方报告促成朝廷或地方的应对记录。 随后可追踪的变化是：赈济、迁徙和损失的实际范围仍须以受灾地区材料分别核验。 可由《英宗实录》正统二年条；《明史·五行志》；受灾府县地方志与奏疏 互相核对；不过，译写候选以编年材料定位；实录记载反映朝廷选择，崇祯期须另以奏疏、地方及敌对政权材料交叉。 因而这里标示的是材料能够支持的行动与制度位置，而不是对动机、地方执行或普通人经验的无保留复原。

\eventanchor{ML-1437-002}{明·1437（正统二年）}
\paragraph{1437（正统二年）：苏北地区发生洪涝灾害} 这一节点应放回本章已经展开的权力、财政、地方社会与边地关系中理解。账本把它出现的条件概括为：自然风险与地方报告促成朝廷或地方的应对记录。 随后可追踪的变化是：赈济、迁徙和损失的实际范围仍须以受灾地区材料分别核验。 可由《英宗实录》正统二年条；《明史·五行志》；受灾府县地方志与奏疏 互相核对；不过，译写候选以编年材料定位；实录记载反映朝廷选择，崇祯期须另以奏疏、地方及敌对政权材料交叉。 因而这里标示的是材料能够支持的行动与制度位置，而不是对动机、地方执行或普通人经验的无保留复原。

\eventanchor{ML-1437-003}{明·1437（正统二年）}
\paragraph{1437（正统二年）：灾情上报、赈恤或河工材料标记本年地方风险，损失与救济效果需回到受灾地区核实。（材料核查重点：诏令或奏议的形成与发受链）} 这一节点应放回本章已经展开的权力、财政、地方社会与边地关系中理解。账本把它出现的条件概括为：自然风险与地方报告促成朝廷或地方的应对记录。 随后可追踪的变化是：赈济、迁徙和损失的实际范围仍须以受灾地区材料分别核验。 可由《英宗实录》正统二年条；《明史·五行志》；受灾府县地方志与奏疏 互相核对；不过，桥接条目只标示可回查的年度问题与材料链；实录有中央选择性，崇祯期尤其需要奏疏、地方、朝鲜及满文材料交叉。；本条特别核对诏令或奏议的形成与发受链。 因而这里标示的是材料能够支持的行动与制度位置，而不是对动机、地方执行或普通人经验的无保留复原。

\eventanchor{ML-1437-004}{明·1437（正统二年）}
\paragraph{1437（正统二年）：灾情上报、赈恤或河工材料标记本年地方风险，损失与救济效果需回到受灾地区核实。（材料核查重点：报称信息的来源、抄传与行政分类）} 这一节点应放回本章已经展开的权力、财政、地方社会与边地关系中理解。账本把它出现的条件概括为：自然风险与地方报告促成朝廷或地方的应对记录。 随后可追踪的变化是：赈济、迁徙和损失的实际范围仍须以受灾地区材料分别核验。 可由《英宗实录》正统二年条；《明史·五行志》；受灾府县地方志与奏疏 互相核对；不过，桥接条目只标示可回查的年度问题与材料链；实录有中央选择性，崇祯期尤其需要奏疏、地方、朝鲜及满文材料交叉。；本条特别核对报称信息的来源、抄传与行政分类。 因而这里标示的是材料能够支持的行动与制度位置，而不是对动机、地方执行或普通人经验的无保留复原。

\eventanchor{MM-1455-ESEN-DEATH}{明·1455（景泰六年）}
\paragraph{1455（景泰六年）：也先在草原内部冲突中被杀，瓦剌汗权继承发生断裂} 这一节点应放回本章已经展开的权力、财政、地方社会与边地关系中理解。账本把它出现的条件概括为：自称大汗后的也先同时面对黄金家族支持者与瓦剌内部首领竞争，个人权威未转化为稳固继承制度。 随后可追踪的变化是：瓦剌的集中领导削弱，明廷北边交涉对象趋于分散；这不等于草原威胁立即消失。 可由《英宗实录》景泰六年北虏条；《明史·瓦剌传》；《蒙古源流》也先死事条 互相核对；不过，死亡及大致年份可由明方边报和蒙古汗统叙事互校；杀害者、月日和各部归属在材料中不一。 因而这里标示的是材料能够支持的行动与制度位置，而不是对动机、地方执行或普通人经验的无保留复原。

\eventanchor{ML-1456-001}{明·1456（景泰七年）}
\paragraph{1456（景泰七年）：湖广苗族叛乱遭镇压} 这一节点应放回本章已经展开的权力、财政、地方社会与边地关系中理解。账本把它出现的条件概括为：继承、官僚协调或皇权运作的压力使问题进入朝廷程序。 随后可追踪的变化是：它影响后续人事与政策议程；动机和派系标签不能由单一叙事裁定。 可由《英宗实录》景泰七年条；《明史·本纪》及相关列传；诏令、奏疏或会典版本 互相核对；不过，译写候选以编年材料定位；实录记载反映朝廷选择，崇祯期须另以奏疏、地方及敌对政权材料交叉。 因而这里标示的是材料能够支持的行动与制度位置，而不是对动机、地方执行或普通人经验的无保留复原。

\eventanchor{ML-1456-002}{明·1456（景泰七年）}
\paragraph{1456（景泰七年）：景泰、天顺时期的继承、人事与诏令链条可在本年材料中定位，政变责任不可由单一史书裁断。} 这一节点应放回本章已经展开的权力、财政、地方社会与边地关系中理解。账本把它出现的条件概括为：继承、官僚协调或皇权运作的压力使问题进入朝廷程序。 随后可追踪的变化是：它影响后续人事与政策议程；动机和派系标签不能由单一叙事裁定。 可由《英宗实录》景泰七年条；《明史·本纪》及相关列传；诏令、奏疏或会典版本 互相核对；不过，桥接条目只标示可回查的年度问题与材料链；实录有中央选择性，崇祯期尤其需要奏疏、地方、朝鲜及满文材料交叉。 因而这里标示的是材料能够支持的行动与制度位置，而不是对动机、地方执行或普通人经验的无保留复原。

\eventanchor{ML-1456-003}{明·1456（景泰七年）}
\paragraph{1456（景泰七年）：军饷、漕运和户籍赋役的本年记录揭示恢复秩序的行政成本，账面额与实收须分列。} 这一节点应放回本章已经展开的权力、财政、地方社会与边地关系中理解。账本把它出现的条件概括为：财政征解、交通工程或货币支付的压力使相关措施进入政务。 随后可追踪的变化是：其法定安排不等于地方执行，后续须以地域材料追踪负担与成效。 可由《英宗实录》景泰七年条；《明会典》相关条目；地方志、黄册/契约/河工材料 互相核对；不过，桥接条目只标示可回查的年度问题与材料链；实录有中央选择性，崇祯期尤其需要奏疏、地方、朝鲜及满文材料交叉。 因而这里标示的是材料能够支持的行动与制度位置，而不是对动机、地方执行或普通人经验的无保留复原。

\eventanchor{ML-1456-004}{明·1456（景泰七年）}
\paragraph{1456（景泰七年）：科举、学校和法律条文在本年制度材料中可追踪，不能据规范文本推定州县实践一致。} 这一节点应放回本章已经展开的权力、财政、地方社会与边地关系中理解。账本把它出现的条件概括为：制度、教育或知识传播的需要推动了相关行动或文本形成。 随后可追踪的变化是：它提供制度和传播史的锚点，不能据此推断社会整体接受。 可由《英宗实录》景泰七年条；《明史·选举志》或《艺文志》；文集、碑刻或书目材料 互相核对；不过，桥接条目只标示可回查的年度问题与材料链；实录有中央选择性，崇祯期尤其需要奏疏、地方、朝鲜及满文材料交叉。 因而这里标示的是材料能够支持的行动与制度位置，而不是对动机、地方执行或普通人经验的无保留复原。

\eventanchor{ML-1457-001}{明·1457（天顺元年）二月11日}
\paragraph{1457（天顺元年）二月11日：前皇帝被军方复职，成为天顺皇帝} 这一节点应放回本章已经展开的权力、财政、地方社会与边地关系中理解。账本把它出现的条件概括为：前期边防、地方武装或跨区域资源调动的压力推动了该行动。 随后可追踪的变化是：它改变了后续调兵、补给或地方秩序的条件；战果和伤亡须分开核对。 可由《英宗实录》天顺元年条；《明史·兵志》及相关列传；边镇、地方志或对方材料 互相核对；不过，译写候选以编年材料定位；实录记载反映朝廷选择，崇祯期须另以奏疏、地方及敌对政权材料交叉。 因而这里标示的是材料能够支持的行动与制度位置，而不是对动机、地方执行或普通人经验的无保留复原。

\subsection{1457年前后的材料线索}

\eventanchor{MM-1457-JINGTAI-DEATH}{明·1457（天顺元年）正月}
\paragraph{1457（天顺元年）正月：英宗复位后景泰帝不久去世，景泰朝的皇位主张随之终结} 这一节点应放回本章已经展开的权力、财政、地方社会与边地关系中理解。账本把它出现的条件概括为：夺门政变使病中的景泰帝失去君位和军政支持，其唯一儿子已早逝，无法以直系储嗣维持景泰朝。 随后可追踪的变化是：英宗一系得以重建继承礼序并处置景泰朝遗臣；景泰帝的陵寝和后世评价遂成为复辟记忆的一部分。 可由《英宗实录》天顺元年正月条；《明史·景帝本纪》；景泰陵碑刻与陵寝调查资料 互相核对；不过，死亡时序和丧葬处置可靠；天顺朝实录以复辟政权立场书写景泰帝病况、名号与责任，不能据此复原全部宫中过程。 因而这里标示的是材料能够支持的行动与制度位置，而不是对动机、地方执行或普通人经验的无保留复原。

\eventanchor{MM-1457-RESTORATION}{明·1457（天顺元年）正月}
\paragraph{1457（天顺元年）正月：石亨、徐有贞、曹吉祥等发动夺门政变，英宗复位} 这一节点应放回本章已经展开的权力、财政、地方社会与边地关系中理解。账本把它出现的条件概括为：景泰帝病重且储位空缺；英宗虽被幽居，仍保有政治象征价值及部分近侍、武臣联系。 随后可追踪的变化是：天顺朝重组中枢人事并以复辟重写景泰朝政治；政变也迅速牵连于谦等守京大臣。 可由《英宗实录》天顺元年正月条；《明史·英宗本纪》；《明史·徐有贞传》 互相核对；不过，政变、复位和参与者大体可靠；谁先谋划、景泰帝病况及各方动机受天顺朝修史塑形。 因而这里标示的是材料能够支持的行动与制度位置，而不是对动机、地方执行或普通人经验的无保留复原。

\eventanchor{MM-1457-YU-QIAN-EXECUTION}{明·1457（天顺元年）正月}
\paragraph{1457（天顺元年）正月：于谦以“谋逆”罪被处死} 这一节点应放回本章已经展开的权力、财政、地方社会与边地关系中理解。账本把它出现的条件概括为：于谦是景泰朝守京和军政的重要大臣，复辟者需要清除与前朝关联最深的政治人物。 随后可追踪的变化是：其死成为天顺朝清洗的标志，也在后世北京保卫战叙事中转化为忠烈象征。 可由《英宗实录》天顺元年正月条；《明史·于谦传》；于谦文集与后世祠祀碑记 互相核对；不过，于谦被杀及所用罪名可靠；谋逆指控的证据链主要出自复辟政权，不能据此断定其有谋反行动。 因而这里标示的是材料能够支持的行动与制度位置，而不是对动机、地方执行或普通人经验的无保留复原。

\eventanchor{ML-1458-001}{明·1458（天顺二年）}
\paragraph{1458（天顺二年）：军饷、漕运和户籍赋役的本年记录揭示恢复秩序的行政成本，账面额与实收须分列。} 这一节点应放回本章已经展开的权力、财政、地方社会与边地关系中理解。账本把它出现的条件概括为：财政征解、交通工程或货币支付的压力使相关措施进入政务。 随后可追踪的变化是：其法定安排不等于地方执行，后续须以地域材料追踪负担与成效。 可由《英宗实录》天顺二年条；《明会典》相关条目；地方志、黄册/契约/河工材料 互相核对；不过，桥接条目只标示可回查的年度问题与材料链；实录有中央选择性，崇祯期尤其需要奏疏、地方、朝鲜及满文材料交叉。 因而这里标示的是材料能够支持的行动与制度位置，而不是对动机、地方执行或普通人经验的无保留复原。

\eventanchor{ML-1458-002}{明·1458（天顺二年）}
\paragraph{1458（天顺二年）：蒙古诸部、朝鲜及西北绿洲的交涉在本年留下多方材料，应以具体使团和地名校正概称。} 这一节点应放回本章已经展开的权力、财政、地方社会与边地关系中理解。账本把它出现的条件概括为：朝贡名义、边境安全与港口贸易的利益使交涉持续发生。 随后可追踪的变化是：它为后续册封、互市或海防安排留下文书与跨政权比较线索。 可由《英宗实录》天顺二年条；《明史·外国传》；《朝鲜王朝实录》、琉球《历代宝案》或海外档案 互相核对；不过，桥接条目只标示可回查的年度问题与材料链；实录有中央选择性，崇祯期尤其需要奏疏、地方、朝鲜及满文材料交叉。 因而这里标示的是材料能够支持的行动与制度位置，而不是对动机、地方执行或普通人经验的无保留复原。
